% Propietario conservado: /Users/ruben/Documents/New project/output/EXTREMA_MEDIA_RAZON_RECIPROCIDAD_ES_EN_20260918/ARTICULO/ES/source/nuclear/sintesis.tex
\chapter{Conclusions sur la conservation, la récupération et la réalisation physique}
\label{x_reciprocidad:nuclear:sintesis}
La structure génératrice conserve les opérations et l’histoire avant sa réalisation numérique. Dans le transport réalisé, cette conservation est formulée comme une identité de formes. Une lecture partielle distribue la charge entre composante visible et complément ; l’inclusion des deux définit une transformation récupérable. Le raffinement prolonge la même identité sans identifier des profondeurs différentes ni remplacer le registre par son intensité.

\section{Conservation, reconstruction et dimension}
L’égalité \(\mathcal U^*(b'\oplus b')\mathcal U=b\oplus b\) réunit lecture et mémoire dans une opération réversible. Pour les compressions à registre séparé, le télescopage conserve le terminal positif et la série complète des compléments. Pour une holonomie unitaire fixe, le terminal est la projection sur son espace de vecteurs fixes ; pour une suite variable, il est déterminé par la limite forte des opérateurs de Gram.

La mémoire dodécaphasique permet de retrouver le registre centré au moyen de deux observations spécifiées, avec une norme d’inverse \(9/4\). L’incidence transporte cette récupération vers le projecteur dimensionnel. Son prolongement extérieur conserve chaque degré au moyen des déterminants de Gram et exige d’inclure les composantes mixtes. Ainsi, les relations entre longueur, aire et volume s’expriment par des applications définies sur le même registre, avec leurs invariants respectifs.

\section{Action et conséquences physiques}
L’action, l’horloge et la base énergétique maintiennent leurs types pendant la composition. Le quotient électronique \(m_ec^2/\hbar\) s’exprime au moyen de la fréquence de l’horloge et du caractère électronique, en conservant le facteur de conversion de base. La normalisation de Maxwell--Dirac élimine les échelles dimensionnelles choisies et laisse le couplage sans dimension. Le transport gravitationnel conserve l’échelle surfacique \(\hbar G/c^3\) ; la solution homogène de rebond maintient ses conditions de matière, de torsion et de constante cosmologique.

La relation de Noether se manifeste lorsque l’observable transportée est une charge de l’action déclarée. Le changement de forme ajoute un terme de connexion au générateur hamiltonien et conserve exactement l’action orientée. L’holonomie d’un retour conserve en outre l’information qui distingue le retour de phase du retour de l’état complet.

\section{Mémoire et spectre quantique}
L’identité \eqref{x_reciprocidad:nuclear:eq:espectro} détermine le spectre symplectique de la covariance réduite par la composante antilinéaire de l’opérateur de mémoire. Pour la boucle
\[
H_1=\begin{pmatrix}2&1\\1&1\end{pmatrix},
\qquad
W_1=\frac19\begin{pmatrix}13&3\\3&10\end{pmatrix},
\]
on obtient \(\det W_1=121/81\), la valeur symplectique \(11/9\) et la pureté \(9/11\). Le spectre de la densité réduite est
\[
\lambda_j=\frac9{10}\,10^{-j},\qquad j=0,1,2,\ldots,
\]
et son entropie vérifie
\[
\frac{S}{k_B}=\frac{10}{9}\log10-\log9.
\]
Les chiffres proviennent du transport et de l’état fixés dans le théorème. L’équivalence spectrale avec une densité de Gibbs attribue une température au Hamiltonien choisi ; cette équivalence conserve sa distinction par rapport à un processus dynamique de thermalisation.

\section{Contenido fundacional}
L’articulation entre nombre, incidence et dynamique réside dans la conservation d’un antécédent commun et dans les opérations qui relient ses réalisations. L’article explicite cette articulation au moyen de reconstructions, de bornes, de spectres et d’actions. La génération des valeurs et le transport des symétries conservent une même provenance ; leurs conséquences observables sont calculées ensuite avec le lecteur et l’état qui interviennent effectivement. Cet ordre fournit une formulation mathématique précise de la mémoire dynamique et de ses réalisations physiques.
